\section{2C与2B的深层分化}

\subsection{2C端的感知停滞}

姚顺雨观察到一个有趣的现象：当人们想到AI Super App时，现在脑中浮现的是两个产品——ChatGPT和Claude Code，分别代表2C和2B的典范。

\begin{warningbox}{2C端的悖论}
对大部分人来说，今天使用ChatGPT的体验与一年前相比变化已经不大。虽然抽象代数、伽罗华理论等高端数学能力显著增强了，但普通用户感受不到。在中国，大部分人仍把AI当作``搜索引擎加强版''，不知如何激发模型的真正智能。
\end{warningbox}

\subsection{2B端的线性价值}

\begin{importantbox}{2B端的价值逻辑}
在2B端，智能越高 = 生产力越高 = 收入越多。这些目标之间是高度aligned的。例如一个OPUS 4.5级别的模型，10个任务可能8--9个直接做对，差的模型只做对5--6个。用户不知道哪些做错了，需要花大量精力监控。因此企业用户愿意花溢价（如每月200美元）使用最好的模型。
\end{importantbox}

2C做AI的困难在于：DAU等产品指标往往与模型智能不相关，甚至可能成反比。而2B端模型越好、收入越高，所有事情都是aligned的。

\subsection{2C端的真正瓶颈：Context而非模型}

\begin{knowledgebox}{一个启发性的例子}
``我今天该吃什么？''——这个问题无论今年、去年还是明年问ChatGPT，效果都不好。因为瓶颈不在于更大的模型或更强的RL，而在于\textbf{额外的context}：天气冷不冷、活动范围、家人的偏好等等。如果能把微信聊天记录等额外输入利用好，反而能带来更多用户价值。
\end{knowledgebox}

\subsection{本章小结}

2C端的瓶颈在context和environment而非模型本身；2B端的价值链高度aligned，是更确定的增长方向。


